% !TeX root = ../../Book2.tex
\section{Calderón–Zygmund 分解}
\label{b2:sec:1802}

一个 \(L^1\) 函数的总质量虽有限，高度却可能任意大。若直接用已有的 \(L^2\) 算子界处理它，高尖峰就会造成困难。Calderón--Zygmund 分解的做法是先选出局部平均过大的区域，把其中的函数压到受控高度；被削下来的部分仍保留在原区域内，但总积分被安排为零。这两部分分别适合不同的估计：高度控制使好部分进入 \(L^2\)，零平均使坏部分经过奇异积分后在远处衰减得更快。

\subsection{好/坏分解}
\label{b2:subsec:180201}

沿用上一章的标准半开二进立方体网格 \(\mathcal D\)。对局部可积函数 \(f\)，记
\[
 f_Q=\frac1{|Q|}\int_Q f(y)\,\mathrm dy.
\]
选择坏区域时要看 \(|f|\) 的平均，避免正负抵消掩盖尖峰；替换函数值时则取 \(f\) 本身的平均。两种平均的作用不同。下面给出\term[calderon-zygmund-fenjie]{Calderón--Zygmund 分解}[Calderón--Zygmund decomposition]的精确形式。

\Needspace{6\baselineskip}
\begin{theorem}[Calderón--Zygmund 分解]\label{b2:thm:18-cz-decomposition}
设 \(d\geq1\) 为整数，\(f\in L^1(\mathbb R^d;\mathbb C)\)，\(\alpha>0\)。存在至多可数个互不相交的半开二进立方体 \(Q_j\)，令 \(E=\bigcup_jQ_j\)，并定义
\[
 g(x)=\begin{cases}
 f(x),&x\notin E,\\
 f_{Q_j},&x\in Q_j,
 \end{cases}
 \quad b_j=(f-f_{Q_j})\cdot\mathbf1_{Q_j}.
\]
则 \(f=g+\sum_jb_j\) 几乎处处，级数在 \(L^1\) 中收敛，并有
\[
 \begin{gathered}
 \alpha<\frac1{|Q_j|}\int_{Q_j}|f|\leq2^d\alpha,
 \quad \sum_j|Q_j|\leq\alpha^{-1}\|f\|_1,\\
 \|g\|_\infty\leq2^d\alpha,
 \quad\|g\|_1\leq\|f\|_1,
 \quad\int b_j=0,
 \quad\sum_j\|b_j\|_1\leq2\|f\|_1.
 \end{gathered}
\]
若另外 \(f\in L^2(\mathbb R^d)\)，则 \(g\in L^2\)，\(\sum_jb_j\) 也在 \(L^2\) 中收敛，并且
\[
 \|g\|_2^2\leq2^d\alpha\|f\|_1,
 \quad \|f\|_2^2=\|g\|_2^2+\sum_j\|b_j\|_2^2.
\]
\end{theorem}
\begin{proof}
用\cref{b2:thm:17-dyadic-level}选择满足
\(|Q|^{-1}\int_Q|f|>\alpha\) 的极大二进立方体。回忆该选择的理由：二进立方体相交时必有一个包含另一个；沿祖先向上，平均值不超过 \(\|f\|_1/|Q|\)，所有足够大的祖先都不再超水平。在仍超水平的有限多个祖先中取最大的一个，就得到所需的极大者。这些极大者互不相交，其二进父立方体 \(\widetilde Q_j\) 的平均不超过 \(\alpha\)。因此
\[
 \int_{Q_j}|f|>\alpha|Q_j|,
 \quad \int_{Q_j}|f|\leq\int_{\widetilde Q_j}|f|
 \leq\alpha|\widetilde Q_j|=2^d\alpha|Q_j|.
\]
对第一式求和得到坏区域的测度界。对 \(E\) 外的点，所有二进平均都不超过 \(\alpha\)；由上一章已经证明的二进微分结论，\(|f|\leq\alpha\) 在该处几乎处处成立。

按陈述定义 \(g\)。在 \(E\) 内，
\[
 |f_{Q_j}|\leq\frac1{|Q_j|}\int_{Q_j}|f|\leq2^d\alpha,
 \quad \int_{Q_j}|g|=|Q_j|\cdot|f_{Q_j}|
                          \leq\int_{Q_j}|f|.
\]
结合 \(E\) 外的结论，得到 \(g\) 的两个范数界。每个坏块满足
\[
 \int b_j=\int_{Q_j}f-|Q_j|\cdot f_{Q_j}=0,
 \quad \|b_j\|_1\leq\int_{Q_j}|f|+|Q_j|\cdot|f_{Q_j}|
                         \leq2\int_{Q_j}|f|.
\]
故 \(\sum_j\|b_j\|_1\leq2\|f\|_1\)，级数在 \(L^1\) 中绝对收敛。不交性还说明每个点至多落在一个坏块中，因而逐点表达与 \(L^1\) 极限相同，给出分解等式。绝对值估计和平均值的线性使这个分解同样适用于复值函数。

若 \(f\in L^2\)，则
\[
 \|g\|_2^2\leq\|g\|_\infty\cdot\|g\|_1
                         \leq2^d\alpha\|f\|_1.
\]
又 \(b=f-g\in L^2\)，且 \(b_j=b\cdot\mathbf1_{Q_j}\)。不交性给出
\(\sum_j\|b_j\|_2^2=\|b\|_2^2\)，所以坏块级数的尾部 \(L^2\) 范数趋于零。更精确地，在每个 \(Q_j\) 上，常数函数 \(f_{Q_j}\) 与 \(b_j\) 正交，因为
\[
 \int_{Q_j}b_j\overline{f_{Q_j}}
 =\overline{f_{Q_j}}\int_{Q_j}b_j=0.
\]
在各块上使用平方范数的勾股恒等式，再与 \(E\) 外的积分合并，便得到陈述中的精确 \(L^2\) 恒等式。
\end{proof}
为什么必须用平均值替换坏块，而不是简单地将它置零？若在 \(Q\) 上用常数 \(a\) 替换，剩余部分的积分为
\(\int_Qf-a|Q|\)。使其为零的唯一选择正是 \(a=f_Q\)。因此，好部分保留了每块的总质量，坏部分只保留块内相对于平均值的变化。若改用 \((|f|)_Q\)，一般就失去了这一抵消。好函数的名称仅指高度受控，并不表示它连续或光滑。

\ProseParagraphBreak

以一维函数 \(f=4\cdot\mathbf1_{[0,1/4)}\) 和高度 \(\alpha=1\) 为例。区间 \([0,1/2)\) 上的平均为 \(2\)，其二进父区间 \([0,1)\) 上的平均恰为 \(1\)；所有更大的祖先平均更小。因此唯一的极大坏区间是 \(Q=[0,1/2)\)，分解为
\[
 g=2\cdot\mathbf1_{[0,1/2)},
 \quad b=2\cdot\mathbf1_{[0,1/4)}
             -2\cdot\mathbf1_{[1/4,1/2)}.
\]
\Cref{b2:fig:18-cz-step-model}显示，削去的正面积由右半块的负面积抵消。这里
\(\|f\|_1=\|g\|_1=\|b\|_1=1\)，而平方范数分别为 \(4,2,2\)，可直接验证定理中的恒等式。坏块并没有变得处处小；它获得的是零平均。

\begin{figure}[htbp]
 \centering
 \begin{tikzpicture}[x=3.6cm,y=0.38cm,font=\small]
  \begin{scope}
   \fill[AnalysisBlueLight] (0,0) rectangle (0.25,4);
   \draw[AnalysisBlue,thick] (0,0)--(0,4)--(0.25,4)--(0.25,0)--(0.72,0);
   \draw[->,AnalysisGray] (-0.03,0)--(0.77,0);
   \draw[->,AnalysisGray] (0,-2.3)--(0,4.7);
   \node[above] at (0.35,5) {原函数 \(f\)};
   \node[left] at (0,4) {\(4\)};
   \foreach \x/\lab in {0/0,0.25/{1/4},0.5/{1/2}} {
    \draw (\x,0.12)--(\x,-0.12);
    \node at (\x,-2.8) {\(\lab\)};
   }
  \end{scope}
  \begin{scope}[xshift=3.8cm]
   \fill[AnalysisBlueLight] (0,0) rectangle (0.5,2);
   \draw[AnalysisBlue,thick] (0,0)--(0,2)--(0.5,2)--(0.5,0)--(0.72,0);
   \draw[->,AnalysisGray] (-0.03,0)--(0.77,0);
   \draw[->,AnalysisGray] (0,-2.3)--(0,4.7);
   \node[above] at (0.35,5) {好部分 \(g\)};
   \node[left] at (0,2) {\(2\)};
   \foreach \x/\lab in {0/0,0.25/{1/4},0.5/{1/2}} {
    \draw (\x,0.12)--(\x,-0.12);
    \node at (\x,-2.8) {\(\lab\)};
   }
  \end{scope}
  \begin{scope}[xshift=7.6cm]
   \fill[AnalysisBlueLight] (0,0) rectangle (0.25,2);
   \fill[AnalysisGrayLight] (0.25,0) rectangle (0.5,-2);
   \draw[AnalysisBlue,thick] (0,0)--(0,2)--(0.25,2)--(0.25,-2)--(0.5,-2)--(0.5,0)--(0.72,0);
   \draw[->,AnalysisGray] (-0.03,0)--(0.77,0);
   \draw[->,AnalysisGray] (0,-2.3)--(0,4.7);
   \node[above] at (0.35,5) {坏部分 \(b\)};
   \node[left] at (0,2) {\(2\)};
   \node[left] at (0,-2) {\(-2\)};
   \foreach \x/\lab in {0/0,0.25/{1/4},0.5/{1/2}} {
    \draw (\x,0.12)--(\x,-0.12);
    \node at (\x,-2.8) {\(\lab\)};
   }
  \end{scope}
 \end{tikzpicture}
 \caption[阶梯函数的好坏分解]{在标准二进网格中，对 \(f=4\cdot\mathbf1_{[0,1/4)}\) 取 \(\alpha=1\)。坏区间为 \([0,1/2)\)，以平均值 \(2\) 替换后得到 \(g\)；剩余的 \(b\) 有等量的正、负面积，故 \(\int b=0\)。图中只画出非零部分附近，三者在所示区间外均为零。}
 \label{b2:fig:18-cz-step-model}
\end{figure}

\subsection{零平均与远场估计}
\label{b2:subsec:180202}

设观察点 \(x\) 离坏块很远。此时 \(K(x,y)\) 随块内变量 \(y\) 的变化较小，可以先把它近似为块中心处的常数 \(K(x,c)\)。这个常数乘上坏块的总积分，恰好为零；真正需要估计的仅是核的变化量。下面的引理把这一思路量化，并保留核差分常数的依赖。

\Needspace{6\baselineskip}
\begin{lemma}\label{b2:thm:18-cancellation-tail}
设 \(d\geq1\) 为整数，\(Q\subset\mathbb R^d\) 是边长 \(\ell>0\)、中心 \(c\) 的轴平行立方体。设 \(b\in L^1(\mathbb R^d;\mathbb C)\) 在 \(Q\) 外几乎处处为零，且 \(\int b=0\)。设复值 Borel 可测核 \(K\) 在 \(x\ne y\) 上取有限值，并存在 \(\delta>0\) 和 \(A_K\geq0\)，使
\[
 |K(x,y)-K(x,c)|
 \leq A_K\frac{|y-c|^\delta}{|x-c|^{d+\delta}}
 \quad\text{当 }x\ne c,\ |x-c|\geq2|y-c|.
\]
令 \(Q^*=8\sqrt d\,Q\) 表示同中心、边长为 \(8\sqrt d\,\ell\) 的立方体。则对 \(x\notin Q^*\)，普通积分 \(\int_QK(x,y)\cdot b(y)\,\mathrm dy\) 绝对收敛，而且
\[
 \int_{\mathbb R^d\setminus Q^*}
 \left|\int_QK(x,y)\cdot b(y)\,\mathrm dy\right|\,\mathrm dx
 \leq C_{d,\delta}A_K\|b\|_1,
\]
其中 \(C_{d,\delta}\) 只依赖 \(d,\delta\)，与 \(Q,b,K\) 无关。
\end{lemma}
\begin{proof}
若 \(x\notin Q^*\)、\(y\in Q\)，则
\[
 |x-c|\geq4\sqrt d\,\ell,
 \quad |y-c|\leq\frac{\sqrt d\,\ell}{2}.
\]
因此差分假设适用。固定这样的 \(x\)，\(K(x,c)\) 有限，核差分在 \(Q\) 上有界，故 \(K(x,\cdot)b\) 绝对可积。零平均给出合法的恒等式
\[
 \int_QK(x,y)\cdot b(y)\,\mathrm dy
 =\int_Q\bigl(K(x,y)-K(x,c)\bigr)\cdot b(y)\,\mathrm dy.
\]
取绝对值，并对非负积分使用 Tonelli，得到所求左端不超过
\[
 A_K\left(\frac{\sqrt d\,\ell}{2}\right)^\delta\|b\|_1
 \int_{|x-c|\geq4\sqrt d\,\ell}|x-c|^{-d-\delta}\,\mathrm dx.
\]
记单位球面的面积为 \(\sigma_{d-1}\)；在 \(d=1\) 时取 \(\sigma_0=2\)。极坐标计算最后一项为
\[
 \sigma_{d-1}\int_{4\sqrt d\,\ell}^\infty r^{-1-\delta}\,\mathrm dr
 =\frac{\sigma_{d-1}}\delta(4\sqrt d\,\ell)^{-\delta}.
\]
这正好消去前面的 \(\ell^\delta\)，给出常数
\(C_{d,\delta}=\sigma_{d-1}/(\delta8^\delta)\) 的一个可用选择。
\end{proof}
这里有两个值得分开的量：坏块的 \(L^1\) 大小决定估计的规模，核变化的快慢决定远场能否积分。若只用大小界 \(|K(x,y)|\lesssim|x-c|^{-d}\)，径向积分会出现 \(\int^\infty\mathrm dr/r\)，无法得到有限上界。减去常数项后，多出的 \(|x-c|^{-\delta}\) 才使它收敛。\(\delta>0\) 正是这一计算所需的条件；尾界对 \(A_K\) 的线性依赖，也使后文能够分别跟踪算子的 \(L^2\) 范数和核的光滑常数。

\ProseParagraphBreak

扩大立方体并不改变坏函数，也不使它产生额外抵消。它只将靠近坏块、不能使用核差分的观察点暂时排除。扩大倍数的具体数值并不重要，只要保证观察距离比块的直径大一个固定倍数；相应的测度代价只依赖维数。

\subsection{弱型估计的分解策略}
\label{b2:subsec:180203}

现在把两个估计合在一起。为清楚看出每项假设的用途，暂设 \(T\) 是 \(L^2(\mathbb R^d)\) 上的有界复线性算子，记 \(B=\|T\|_{2\to2}\)。还假设它在每个紧支集 \(L^2\) 输入的支集外，几乎处处由核 \(K\) 的普通积分表示；核对每个立方体中心都满足上一引理的差分条件，且使用共同的 \(\delta,A_K\)。下一节会把这些条件组织成标准算子定义。

\ProseParagraphBreak

给定 \(f\in L^1\cap L^2\) 和输出阈值 \(t>0\)，在高度 \(\alpha=t\) 分解 \(f=g+b\)。好部分使用 \(L^2\) Chebyshev 不等式：
\[
 \begin{aligned}
 |\{|Tg|>t/2\}|
 &\leq\frac4{t^2}\|Tg\|_2^2
 \leq\frac{4B^2}{t^2}\|g\|_2^2\\
 &\leq\frac{2^{d+2}B^2}{t}\|f\|_1.
 \end{aligned}
\]
好部分的有界性把它的平方积分控制在 \(t\|f\|_1\) 的量级，这正是分解带来的收益。

\ProseParagraphBreak

对坏部分，先丢弃扩大坏区域
\[
 E^*=\bigcup_jQ_j^*,
 \quad |E^*|\leq(8\sqrt d)^d\sum_j|Q_j|
          \leq\frac{(8\sqrt d)^d}{t}\|f\|_1.
\]
原来的 \(Q_j\) 互不相交，扩大的 \(Q_j^*\) 可能重叠；这里用的是并集测度不超过测度之和，不需要扩大后仍然不交。

\ProseParagraphBreak

在 \((E^*)^c\) 上使用坏块的零平均。先取有限和 \(b^{(N)}=\sum_{j=1}^Nb_j\)。每个 \(b_j\) 属于 \(L^2\) 且支撑于一个有界立方体，因而支集外核表示与远场引理给出
\[
 \begin{aligned}
 \int_{(E^*)^c}|Tb^{(N)}(x)|\,\mathrm dx
 &\leq\sum_{j=1}^N\int_{(Q_j^*)^c}|Tb_j(x)|\,\mathrm dx\\
 &\leq C_{d,\delta}A_K\sum_{j=1}^N\|b_j\|_1
 \leq2C_{d,\delta}A_K\|f\|_1.
 \end{aligned}
\]
若坏块有无穷多个，定理保证 \(b^{(N)}\to b\) 于 \(L^2\)，于是 \(Tb^{(N)}\to Tb\) 于 \(L^2\)。抽出几乎处处收敛的子列，再对 \((E^*)^c\) 使用 Fatou 引理，得到
\[
 \int_{(E^*)^c}|Tb(x)|\,\mathrm dx
 \leq2C_{d,\delta}A_K\|f\|_1.
\]
这不要求 \(\sum_jTb_j\) 在全空间的 \(L^1\) 中收敛。我们得到的仅是远离坏区域后的积分控制，正好足够用于 Chebyshev 不等式：
\[
 |\{x\notin E^*:|Tb(x)|>t/2\}|
 \leq\frac{4C_{d,\delta}A_K}{t}\|f\|_1.
\]

由 \(Tf=Tg+Tb\)，其 \(t\) 超水平集包含在好部分的 \(t/2\) 超水平集、扩大坏区域，以及坏部分在其外的 \(t/2\) 超水平集三者之并中。合并以上估计即得
\[
 |\{|Tf|>t\}|
 \leq\frac{C_{d,\delta}(1+B^2+A_K)}t\|f\|_1
 \quad(f\in L^1\cap L^2,\ t>0),
\]
其中最后的 \(C_{d,\delta}\) 可以增大以吸收各项维数常数。好部分靠 \(L^2\) 有界性，坏区域靠测度控制，坏部分的远场靠零平均与核差分；三条估计承担着不同任务。由稠密性把算子延拓到全部 \(L^1\)，以及由弱端点推出 \(L^p\) 强型的步骤，将在\secref{b2:sec:1804}完成。
